\documentclass[10pt,a4paper]{amsart}
\usepackage[margin=19mm]{geometry}
\usepackage{amsmath,amssymb,mathtools}
\usepackage[scr=rsfs,cal=euler]{mathalfa}
\usepackage{xfrac}
\usepackage[unicode,hidelinks,bookmarks=false]{hyperref}
\newcommand{\ie}{i.e.\ }
\newcommand{\cf}{cf.\ }
\newcommand{\resp}{respectively\ }
\newcommand{\Subsection}{\subsection}
\providecommand{\nobreakdash}{\nobreak\hbox{-}}
\setlength{\emergencystretch}{2em}
\multlinegap0pt
\usepackage{bm}
\usepackage{bbm}
\usepackage{dsfont}
\usepackage{xspace}
\usepackage{cancel}
\usepackage{mathtools}
\usepackage{amsthm}
\makeatletter
\@ifundefined{Theorem}{\newtheorem{Theorem}{Theorem}}{}
\@ifundefined{Definition}{\newtheorem{Definition}[Theorem]{Definition}}{}
\@ifundefined{Lemma}{\newtheorem{Lemma}[Theorem]{Lemma}}{}
\@ifundefined{Proposition}{\newtheorem{Proposition}[Theorem]{Proposition}}{}
\makeatother
\DeclareMathOperator{\Par}{\mathbb{Y}}
\DeclareMathOperator{\SYT}{\mathrm{SYT}}
\DeclareMathOperator{\sH}{\mathscr{H}}
\title[Murnaghan-Type Representations for the Positive Elliptic Hal]{Murnaghan-Type Representations for the Positive Elliptic Hall Algebra}
\author{Milo Bechtloff Weising}
\date{Source: arXiv:2405.00756v2 (CC BY 4.0)}
\begin{document}
\maketitle

\noindent\emph{Source and licence.} Murnaghan-Type Representations for the Positive Elliptic Hall Algebra. Version arXiv:2405.00756v2 (CC BY 4.0), distributed under the Creative Commons Attribution 4.0 International licence, as recorded in the arXiv licence metadata for this exact version and shown on the abstract page https://arxiv.org/abs/2405.00756v2. Full rights notice: licenses/arxiv-2405.00756-CC-BY-4.0.txt.\par\smallskip

\noindent\textit{Editorial scope.} Counted unit: the Proposition whose source label is \texttt{relations for connecting maps} in the article (source0.tex lines 1352--1363, source label \texttt{relations for connecting maps}), delivered together with the article's own complete original proof of it (source0.tex lines 1365--1451, one \texttt{proof} environment of 87 lines). No mathematics is written, repaired or re-derived: statement and proof are the article's own text line for line. Required context retained uncounted: finite hecke map (lines 644--654); action of T inverse (lines 1093--1102); connecting maps def (lines 1342--1348). Every definition, display and label of those lines is retained unchanged. Omitted from this package and not redistributed: 1--643, 655--1092, 1103--1341, 1349--1351, 1364--1364, 1452--2600. Nothing inside a retained excerpt is omitted or altered: every retained body is delivered exactly as the article's lines are written, so no omission receipt of any kind accompanies this package. Citations: the retained excerpts carry no citation command at all, so no bibliography entry of the article is transcribed and no reference is redistributed. Reference adaptation: none was needed and none was made; every reference inside the delivered unit resolves inside it, so no unresolved reference or citation is produced. Printed slips kept literal: no printed word, symbol, hypothesis or number of the counted statement or of its proof was altered. Layout and class: the article's own class is \texttt{sigma} with packages including ifpdf, ytableau, mathrsfs, bbm; the wrapper is the lane's pinned \texttt{amsart} editorial wrapper at 10pt on A4, which restates the theorem-like environments the retained text needs and adds the article's own macro definitions that the retained text needs (\texttt{\textbackslash Par}, \texttt{\textbackslash SYT}, \texttt{\textbackslash sH}), with extra wrapper packages (bm, bbm, dsfont, xspace, cancel, mathtools). The delivered numbering is the unit's own. Assets: no retained span contains a figure environment, a graphics-inclusion command, a PDF-page inclusion, a table or a TikZ picture; no image file is copied into this package, the package declares no asset dependency and no image file is redistributed with it. Licence: version arXiv:2405.00756v2 of this article is distributed under the Creative Commons Attribution 4.0 International licence, as recorded in the arXiv licence metadata for this exact version and shown on the abstract page https://arxiv.org/abs/2405.00756v2. A full rights notice is delivered at licenses/arxiv-2405.00756-CC-BY-4.0.txt. Characters: every retained excerpt is pure ASCII, so the delivered engine, XeLaTeX with its default Latin Modern text font, reports no missing character.
\begin{Lemma}\label{finite hecke map}
 Let $\lambda \in \Par$ and $n \geq n_{\lambda}$. Let $\square_0$ be the unique square in the skew-diagram \smash{$\lambda^{(n+1)} / \lambda^{(n)}$}. Consider the map \smash{$\mathfrak{q}_{\lambda}^{(n)}\colon\lambda^{(n+1)}\rightarrow \lambda^{(n)}$} given for \smash{$\tau \in \SYT\bigl(\lambda^{(n+1)}\bigr)$} as
 \[
 \mathfrak{q}_{\lambda}^{(n)}(e_{\tau}) := \begin{cases}
 e_{\tau|_{\lambda^{(n)}}}, & \tau(\square_0) = n+1,\\
 0, & \tau(\square_0) \neq n+1.
 \end{cases}
 \]

 Then $\mathfrak{q}_{\lambda}^{(n)}$ is a~$\sH_{n}$-module map.
\end{Lemma}

\begin{Lemma}\label{action of T inverse}
 For $1\leq i \leq n-1$ and $a \geq 0$,
 \[
 \bigl(tT_i^{-1}\bigr)X_{i+1}^{a} = X_i^{a}\bigl(tT_i^{-1}\bigr) + (t-1)X_{i+1}\frac{X_i^{a}-X_{i+1}^{a}}{X_i-X_{i+1}}.
 \]
 Further, every monomial occurring in the term \smash{$X_{i+1}\frac{X_i^{a}-X_{i+1}^{a}}{X_i-X_{i+1}}$} is strictly lower than $X_i^{a}$ with respect to the Bruhat ordering $\preceq$. Consequently, it follows that for any $\alpha \in \mathbb{Z}^n_{\geq 0}$ with $s_i(\alpha) \succeq \alpha$ the following expansion holds for some scalars $c_{\beta}$:
 \[
 \bigl(tT_i^{-1}\bigr)X^{\alpha} = X^{s_i(\alpha)}\bigl(tT_i^{-1}\bigr) + \sum_{\beta \prec s_i(\alpha)} c_{\beta}X^{\beta}.
 \]
\end{Lemma}

\begin{Definition}\label{connecting maps def}
 Let $\lambda \in \Par$. For $n \geq n_{\lambda}$, define \smash{$\Phi^{(n)}_{\lambda}\colon V_{\lambda^{(n+1)}} \rightarrow V_{\lambda^{(n)}}$} as the $\mathbb{Q}(q,t)$-linear map given on any element $X^{\alpha}\otimes v \in V_{\lambda^{(n+1)}}$ by
 \[
 \Phi^{(n)}_{\lambda}(X^{\alpha}\otimes v) = \mathbbm{1}(\alpha_{n+1} = 0) X_1^{\alpha_1}\cdots X_n^{\alpha_n}\otimes \mathfrak{q}^{(n)}_{\lambda}(v).
 \]

\end{Definition}

\begin{Proposition}\label{relations for connecting maps}
 The map \smash{$\Phi^{(n)}_{\lambda}$} satisfies the following relations:
 \begin{itemize}
\itemsep=0pt
 \item \smash{$\Phi^{(n)}_{\lambda}T_i = T_i\Phi^{(n)}_{\lambda}$} for $1\leq i \leq n-1$.
 \item \smash{$\Phi^{(n)}_{\lambda}X_i = X_i\Phi^{(n)}_{\lambda}$} for $1\leq i \leq n$.
 \item \smash{$\Phi^{(n)}_{\lambda}X_{n+1} = 0$}.
 \item \smash{$\Phi^{(n)}_{\lambda} t^{-n}\pi_{n+1}T_n = t^{-(n-1)}\pi_n \Phi^{(n)}_{\lambda}$}.
 \item \smash{$\Phi^{(n)}_{\lambda}\theta_i^{(n+1)} = \theta_i^{(n)}\Phi^{(n)}_{\lambda}$} for $1\leq i \leq n$.
 \item \smash{$\Phi^{(n)}_{\lambda}\bigl(\theta_{n+1}^{(n+1)} - t^{n-|\lambda|}\bigr) = 0$}.
 \end{itemize}
\end{Proposition}

\begin{proof}
 From Lemma~\ref{finite hecke map} and Definition~\ref{connecting maps def}, it follows immediately for all $1\leq i \leq n-1$ and~${1\leq j \leq n}$, that \smash{$\Phi^{(n)}_{\lambda}T_i = T_i\Phi^{(n)}_{\lambda}$}, \smash{$\Phi^{(n)}_{\lambda}X_j = X_j\Phi^{(n)}_{\lambda}$}, and \smash{$\Phi^{(n)}_{\lambda}X_{n+1} = 0$}.

 Let \smash{$X^{\alpha}\otimes v \in V_{\lambda^{(n+1)}}$}.
 By direct calculation, we find
 \begin{gather*}
 \Phi^{(n)}_{\lambda}t^{-n}\pi_{n+1}T_n\bigl(X_1^{\alpha_1}\cdots X_{n+1}^{\alpha_{n+1}}\otimes v\bigr) \\
 \qquad = \Phi^{(n)}_{\lambda}t^{-n}\pi_{n+1} X_1^{\alpha_1}\cdots X_{n-1}^{\alpha_{n-1}}T_n\bigl(X_n^{\alpha_n}X_{n+1}^{\alpha_{n+1}}\otimes v\bigr) \\
 \qquad = \Phi^{(n)}_{\lambda}t^{-n}X_2^{\alpha_1}\cdots X_n^{\alpha_{n-1}}\pi_{n+1}T_n\bigl(X_n^{\alpha_n}X_{n+1}^{\alpha_{n+1}} \otimes v\bigr) \\
 \qquad = t^{-n}X_2^{\alpha_1}\cdots X_n^{\alpha_{n-1}} \Phi^{(n)}_{\lambda} \pi_{n+1} T_n\bigl(X_n^{\alpha_n}X_{n+1}^{\alpha_{n+1}} \otimes v\bigr) \\
 \qquad = t^{-n}X_2^{\alpha_1}\cdots X_n^{\alpha_{n-1}} \\
 \phantom{ \qquad =}{} \times \Phi^{(n)}_{\lambda} \pi_{n+1} \left(X_{n+1}^{\alpha_n}X_n^{\alpha_{n+1}}T_n\otimes v + (1-t)X_n \frac{X_n^{\alpha_n}X_{n+1}^{\alpha_{n+1}}-X_{n+1}^{\alpha_n}X_{n}^{\alpha_{n+1}}}{X_n -X_{n+1}}\otimes v \right) \\
 \qquad = t^{-n}X_2^{\alpha_1}\cdots X_n^{\alpha_{n-1}} \Phi^{(n)}_{\lambda} \left( q^{\alpha_n}X_1^{\alpha_n}X_{n+1}^{\alpha_{n+1}} \pi_{n+1}T_n \otimes v \right. \\
\phantom{ \qquad =}{} \left. + (1-t)X_{n+1}\pi_{n+1}\frac{X_n^{\alpha_n}X_{n+1}^{\alpha_{n+1}}-X_{n+1}^{\alpha_n}X_{n}^{\alpha_{n+1}}}{X_n -X_{n+1}} \otimes v \right) \\
 \qquad = \mathbbm{1}(\alpha_{n+1}=0)t^{-n}q^{\alpha_n}X_1^{\alpha_n}X_2^{\alpha_1}\cdots X_n^{\alpha_{n-1}}\Phi^{(n)}_{\lambda}(1 \otimes \rho_{n+1}(\pi_{n+1}T_n)v) \\
 \qquad = \mathbbm{1}(\alpha_{n+1}=0)t^{-n}q^{\alpha_n}X_1^{\alpha_n}X_2^{\alpha_1}\cdots X_n^{\alpha_{n-1}}\Phi^{(n)}_{\lambda}\bigl(1 \otimes t^nT_1^{-1}\cdots T_{n-1}^{-1}v\bigr) \\
 \qquad = \mathbbm{1}(\alpha_{n+1}=0)q^{\alpha_n}X_1^{\alpha_n}X_2^{\alpha_1}\cdots X_n^{\alpha_{n-1}} \otimes T_1^{-1}\cdots T_{n-1}^{-1}\mathfrak{q}_{\lambda}^{(n)}(v).
\end{gather*}

On the other hand, we see
\begin{gather*}
 t^{-(n-1)}\pi_n \Phi^{(n)}_{\lambda}\bigl(X_1^{\alpha}\cdots X_{n+1}^{\alpha_{n+1}}\otimes v\bigr) \\
 \qquad =\mathbbm{1}(\alpha_{n+1}=0) t^{-(n-1)}\pi_n\bigl(X_1^{\alpha_1}\cdots X_n^{\alpha_n}\bigr) \otimes \mathfrak{q}_{\lambda}^{(n)}(v) \\
 \qquad = \mathbbm{1}(\alpha_{n+1} =0) t^{-(n-1)}q^{\alpha_n}X_1^{\alpha_n}X_2^{\alpha_1}\cdots X_n^{\alpha_{n-1}}\otimes \rho_n(\pi_n)\bigl(\mathfrak{q}_{\lambda}^{(n)}(v)\bigr) \\
 \qquad = \mathbbm{1}(\alpha_{n+1} =0)q^{\alpha_n}X_1^{\alpha_n}X_2^{\alpha_1}\cdots X_n^{\alpha_{n-1}}\otimes T_1^{-1}\cdots T_{n-1}^{-1}\mathfrak{q}_{\lambda}^{(n)}(v).
\end{gather*}

Therefore, \smash{$\Phi^{(n)}_{\lambda} t^{-n}\pi_{n+1}T_n = t^{-(n-1)}\pi_n \Phi^{(n)}_{\lambda}$} as desired.

Now, let $1\leq i \leq n$. We see that
\begin{align*}
 \Phi^{(n)}_{\lambda} \theta_i^{(n+1)}&= \Phi^{(n)}_{\lambda} t^{-(n-i+1)}T_{i-1}^{-1}\cdots T_1^{-1} \pi_{n+1} T_n\cdots T_i \\
 & = t^{i-1}T_{i-1}^{-1}\cdots T_1^{-1}\bigl(\Phi^{(n)}_{\lambda}t^{-n}\pi_{n+1} T_n\bigr) T_{n-1}\cdots T_i \\
 & = t^{i-1}T_{i-1}^{-1}\cdots T_1^{-1} \bigl(t^{-(n-1)}\pi_n \Phi^{(n)}_{\lambda}\bigr)T_{n-1}\cdots T_i \\
 & = t^{-(n-i)} T_{i-1}^{-1}\cdots T_1^{-1} \pi_n T_{n-1}\cdots T_i\Phi^{(n)}_{\lambda}
 = \theta_i^{(n)}\Phi^{(n)}_{\lambda}.
\end{align*}

Now, let \smash{$\alpha \in \mathbb{Z}^{n+1}_{\geq 0}$} and $\tau \in \SYT\bigl(\lambda^{(n+1)}\bigr)$. We find
\begin{align*}
 \Phi^{(n)}_{\lambda} \theta_{n+1}^{(n+1)}(X^{\alpha} \otimes e_{\tau}) &
 = \Phi^{(n)}_{\lambda}T_n^{-1}\cdots T_1^{-1} \pi_{n+1}(X^{\alpha} \otimes e_{\tau}) \\
 & = \Phi^{(n)}_{\lambda}T_n^{-1}\cdots T_1^{-1} q^{\alpha_{n+1}}X_1^{\alpha_{n+1}}X_2^{\alpha_1}\cdots X_{n+1}^{\alpha_n} \otimes \rho_{n+1}(\pi_{n+1})e_{\tau} \\
 & = q^{\alpha_{n+1}}\Phi^{(n)}_{\lambda}T_n^{-1}\cdots T_1^{-1}X_1^{\alpha_{n+1}}X_2^{\alpha_1}\cdots X_{n+1}^{\alpha_n} \bigl(1\otimes t^nT_1^{-1}\cdots T_n^{-1}(e_{\tau})\bigr).
\end{align*}
If $\alpha_{n+1} >0$, then this evaluates to $0$ since
\[
\Phi^{(n)}_{\lambda}T_n^{-1}\cdots T_1^{-1}X_1 = t^{-n}\Phi^{(n)}_{\lambda} X_{n+1}T_n \cdots T_1 =0.
\]
Hence,
\[
 \Phi^{(n)}_{\lambda} \theta_{n+1}^{(n+1)}(X^{\alpha} \otimes e_{\tau}) = \mathbbm{1}(\alpha_{n+1} =0)\Phi^{(n)}_{\lambda}T_n^{-1}\cdots T_1^{-1}X_2^{\alpha_1}\cdots X_{n+1}^{\alpha_n} \bigl(1\otimes t^nT_1^{-1}\cdots T_n^{-1}(e_{\tau})\bigr).
\]

By repeatedly applying Lemma~\ref{action of T inverse}, we see that as maps $V_{\lambda^{(n+1)}} \rightarrow V_{\lambda^{(n)}}$
\begin{gather*}
 \Phi^{(n)}_{\lambda}T_n^{-1}\cdots T_2^{-1} \bigl(T_{1}^{-1}X_2^{\alpha_1}\bigr) X_3^{\alpha_2}\cdots X_{n+1}^{\alpha_n} \\
 \qquad = \Phi^{(n)}_{\lambda}T_n^{-1}\cdots T_2^{-1}\left(X_{1}^{\alpha_1}T_1^{-1} + \bigl(1-t^{-1}\bigr)X_2 \frac{X_1^{\alpha_1}-X_2^{\alpha_1}}{X_1-X_2}\right)X_3^{\alpha_2}\cdots X_{n+1}^{\alpha_n} \\
 \qquad =X_{1}^{\alpha_1}\Phi^{(n)}_{\lambda}T_n^{-1}\cdots T_1^{-1}X_3^{\alpha_2}\cdots X_{n+1}^{\alpha_n} \\
 \phantom{\qquad =}{} + \bigl(1-t^{-1}\bigr)\Phi^{(n)}_{\lambda}T_n^{-1}\cdots T_2^{-1}X_2 \frac{X_1^{\alpha_1}-X_2^{\alpha_1}}{X_1-X_2}X_3^{\alpha_2}\cdots X_{n+1}^{\alpha_n} \\
 \qquad = X_{1}^{\alpha_1}\Phi^{(n)}_{\lambda}T_n^{-1}\cdots T_1^{-1}X_3^{\alpha_2}\cdots X_{n+1}^{\alpha_n} \\
 \phantom{\qquad =}{} + \bigl(1-t^{-1}\bigr)t^{-(n-2)}\Phi^{(n)}_{\lambda} X_{n+1} T_{n-1}\cdots T_2 \frac{X_1^{\alpha_1}-X_2^{\alpha_1}}{X_1-X_2}X_3^{\alpha_2}\cdots X_{n+1}^{\alpha_n} \\
 \qquad = X_{1}^{\alpha_1}\Phi^{(n)}_{\lambda}T_n^{-1}\cdots T_1^{-1}X_3^{\alpha_2}\cdots X_{n+1}^{\alpha_n} + 0 \\
 \qquad = X_{1}^{\alpha_1}\Phi^{(n)}_{\lambda}T_n^{-1}\cdots T_3^{-1}(T_2^{-1} X_3^{\alpha_2})T_1^{-1}X_4^{\alpha_3}\cdots X_{n+1}^{\alpha_n} \\
 \qquad = \cdots \\
 \qquad = X_{1}^{\alpha_1}X_2^{\alpha_2}\Phi^{(n)}_{\lambda}T_n^{-1}\cdots T_1^{-1}X_4^{\alpha_3}\cdots X_{n+1}^{\alpha_n} \\
 \qquad = \cdots \\
 \qquad = X_1^{\alpha_1}\cdots X_{n}^{\alpha_{n}}\Phi^{(n)}_{\lambda}T_n^{-1}\cdots T_1^{-1}.
\end{gather*}

As usual, let $\square_0$ denote the unique square of the skew diagram $\lambda^{(n+1)}/\lambda^{(n)}$. Returning to our main calculation now shows
\begin{gather*}
 \Phi^{(n)}_{\lambda} \theta_{n+1}^{(n+1)}(X^{\alpha} \otimes e_{\tau}) \\
 \qquad =\mathbbm{1}(\alpha_{n+1} =0)\Phi^{(n)}_{\lambda}T_n^{-1}\cdots T_1^{-1}X_2^{\alpha_1}\cdots X_{n+1}^{\alpha_n} \bigl(1\otimes t^nT_1^{-1}\cdots T_n^{-1}(e_{\tau})\bigr) \\
 \qquad = \mathbbm{1}(\alpha_{n+1} =0)X_1^{\alpha_1}\cdots X_{n}^{\alpha_{n}}\Phi^{(n)}_{\lambda}T_n^{-1}\cdots T_1^{-1} \bigl(1\otimes t^nT_1^{-1}\cdots T_n^{-1}(e_{\tau})\bigr) \\
 \qquad = \mathbbm{1}(\alpha_{n+1} =0)X_1^{\alpha_1}\cdots X_{n}^{\alpha_{n}}\Phi^{(n)}_{\lambda}\bigl(1 \otimes t^nT_n^{-1}\cdots T_1^{-1}T_1^{-1}\cdots T_n^{-1}(e_{\tau})\bigr) \\
 \qquad = \mathbbm{1}(\alpha_{n+1} =0)X_1^{\alpha_1}\cdots X_{n}^{\alpha_{n}}\Phi^{(n)}_{\lambda}\bigl(1 \otimes \overline{\theta}_{n+1}^{(n+1)}(e_{\tau})\bigr) \\
 \qquad = \mathbbm{1}(\alpha_{n+1} =0)X_1^{\alpha_1}\cdots X_{n}^{\alpha_{n}}\Phi^{(n)}_{\lambda}\bigl(1 \otimes t^{c_{\tau}(n+1)}e_{\tau}\bigr) \\
 \qquad = t^{c_{\tau}(n+1)}\mathbbm{1}(\alpha_{n+1} =0)X_1^{\alpha_1}\cdots X_{n}^{\alpha_{n}}\otimes \mathfrak{q}_{\lambda}^{(n)}(e_{\tau}) \\
 \qquad = t^{c(\square_0)}\mathbbm{1}(\alpha_{n+1} =0)\mathbbm{1}( \tau(\square_0) = n+1)X_1^{\alpha_1}\cdots X_{n}^{\alpha_{n}}\otimes e_{\tau|{\lambda^{(n)}}} \\
 \qquad = t^{n-|\lambda|} \mathbbm{1}(\alpha_{n+1} =0)\mathbbm{1}( \tau(\square_0) = n+1)X_1^{\alpha_1}\cdots X_{n}^{\alpha_{n}}\otimes e_{\tau|{\lambda^{(n)}}} = \Phi^{(n)}_{\lambda}\bigl(t^{n-|\lambda|}X^{\alpha}\otimes e_{\tau}\bigr).
\end{gather*}

Therefore,
\smash{$
\Phi^{(n)}_{\lambda}\bigl( \theta_{n+1}^{(n+1)} - t^{n-|\lambda|}\bigr) = 0$}.
\end{proof}

\end{document}
